\documentclass[11pt,a4paper]{article}
\usepackage[margin=2.2cm]{geometry}
\usepackage[T1]{fontenc}
\usepackage[utf8]{inputenc}
\usepackage{lmodern}
\usepackage{booktabs}
\usepackage{longtable}
\usepackage{tabularx}
\usepackage{enumitem}
\usepackage{fancyhdr}
\usepackage{titlesec}
\usepackage{hyperref}
\usepackage{xcolor}
\usepackage{setspace}
\usepackage{microtype}

\hypersetup{
  colorlinks=true,
  linkcolor=blue!50!black,
  urlcolor=blue!60!black,
  pdftitle={Open Clinical Record --- User Manual},
  pdfauthor={Open Clinical Record Internship Project}
}

\pagestyle{fancy}
\fancyhf{}
\fancyhead[L]{\small Open Clinical Record}
\fancyhead[R]{\small User Manual}
\fancyfoot[C]{\thepage}
\renewcommand{\headrulewidth}{0.4pt}

\titleformat{\section}{\large\bfseries\color{blue!40!black}}{\thesection}{0.8em}{}
\titleformat{\subsection}{\normalsize\bfseries}{\thesubsection}{0.7em}{}
\setlist[itemize]{leftmargin=1.4em,itemsep=0.25em}
\setlist[enumerate]{leftmargin=1.6em,itemsep=0.25em}
\onehalfspacing

\begin{document}

\begin{titlepage}
\centering
\vspace*{1.5cm}
{\LARGE\bfseries Open Clinical Record}\\[0.4cm]
{\Large User Manual}\\[1.0cm]
{\large Outpatient Electronic Medical Record}\\[0.4cm]
{\large Version 1.0 --- Week 8 Finalization}\\[0.4cm]
{\large 27 September 2026}\\[1.5cm]
\begin{tabular}{ll}
\toprule
\textbf{Audience} & Receptionists, nurses, doctors, administrators \\
\textbf{Product} & Open Clinical Record (OCR) \\
\textbf{Companion docs} & SRS v2.0, Technical Documentation \\
\bottomrule
\end{tabular}
\vfill
{\small Treat all patient information as confidential. Sign out on shared workstations.}
\end{titlepage}

\tableofcontents
\newpage

\section{Introduction}

\subsection{What this system is for}
Open Clinical Record helps a clinic run three connected workflows:
\begin{enumerate}
  \item \textbf{Register and find patients}
  \item \textbf{Book and manage appointments}, including check-in
  \item \textbf{Document clinical visits} so history accumulates safely over time
\end{enumerate}
It is deliberately focused. It is not a full hospital system (no pharmacy, laboratory orders, or billing modules in this version).

\subsection{Who uses it}
\begin{tabularx}{\textwidth}{@{}p{3.6cm}X@{}}
\toprule
\textbf{Role} & \textbf{Typical work} \\
\midrule
Receptionist & Register patients, book/cancel/reschedule, check-in, mark deceased \\
Nurse & Support the queue, record vitals and notes on visits \\
Doctor & Review the chart, document and finalize visits, mark/clear deceased \\
Administrator & Users and audit; system oversight \\
\bottomrule
\end{tabularx}

Your menu only shows what your role is allowed to use. The server still blocks forbidden actions.

\section{Getting Started}

\subsection{Sign in}
\begin{enumerate}
  \item Open the application (local development: \texttt{http://localhost:5173}).
  \item Enter your email and password.
  \item You arrive on a role-aware dashboard.
\end{enumerate}

\subsection{Training accounts (development only)}
\begin{tabularx}{\textwidth}{@{}X X X@{}}
\toprule
\textbf{Email} & \textbf{Role} & \textbf{Default password} \\
\midrule
\texttt{desk@clinic.local} & Receptionist & \texttt{Dev@12345} \\
\texttt{nurse@clinic.local} & Nurse & \texttt{Dev@12345} \\
\texttt{doctor@clinic.local} & Doctor & \texttt{Dev@12345} \\
\texttt{admin@clinic.local} & Admin & \texttt{Dev@12345} \\
\bottomrule
\end{tabularx}

\vspace{0.4em}
\noindent\textit{Never use these passwords outside training. Change them before production-like deployment.}

\subsection{Navigation}
Sidebar is role-aware: Dashboard, Patients, Register, Chart, Appointments, Check-in/Queue, Profile, Admin (Admin only).

\section{Receptionist Guide}

\subsection{Register a patient}
Open \textbf{Register patient}, search first to avoid duplicates, enter demographics, save. A medical record number is assigned.

\subsection{Book an appointment}
Open \textbf{New appointment}, select patient, date, duration, slot, type, then create. Cannot book for Deceased patients.

\subsection{Check-in}
Open \textbf{Check-in / Queue}, check in the patient. Status becomes Checked in; a Draft visit is created. Completed/Cancelled/No-show stay out of the active queue.

\subsection{Cancel and reschedule}
Cancel requires a non-empty reason. Reschedule is allowed for Scheduled, Waiting, Cancelled, No-show; blocked for CheckedIn, InProgress, Completed.

\subsection{Mark deceased}
Receptionists may mark deceased. Clearing is Doctor/Admin only.

\section{Nurse Guide}
Open queue and chart. On the Draft visit enter vitals and notes, then save. Nurses cannot mark deceased.

\section{Doctor Guide}
Review chart history. Document the Draft visit; set \textbf{Final} when complete. Final visits are immutable; returns create a new visit. Finalizing completes the linked appointment for the queue. Doctors may mark and clear deceased.

\section{Administrator Guide}
Manage users and review audit events. Admin is not a default clinical author.

\section{Statuses}
\textbf{Appointment:} Scheduled, Waiting, Checked in, In progress, Completed, Cancelled, No-show.\\
\textbf{Visit:} Draft (editable), Final (immutable), Cancelled.

\section{Troubleshooting}
Blank page: hard refresh, confirm API, re-login. Permission denied: wrong role. Cancel blocked: enter reason. Missing from queue: terminal status or date filter.

\section{Privacy}
Access only needed records. Sign out on shared workstations. Use synthetic data in screenshots.

\section*{Document information}
User Manual v1.0 (27 September 2026). See SRS v2.0 and Technical Documentation for requirements and architecture.

\end{document}
